% Section I: Introduction - from standard QFT to topological organization
% Purpose: State the physics problem, not an abstract definition
% Target audience: Graduate student familiar with Schwartz-level QFT
%
% Key questions to answer:
% (i) what is topological structure inside ordinary QFT?
% (ii) why does Chern-Simons theory isolate that structure so cleanly?
% (iii) how does this lead to the formal TQFT framework?
%
% Status: TO BE WRITTEN (write last, after main content is complete)
% Sources: None yet - this frames the entire paper
%
% Writing notes:
% - Do NOT open with abstract TQFT definition
% - Open with physics problem that Schwartz-trained student recognizes
% - State early that paper will answer the three questions above
% - Keep to 2-3 pages maximum
% - Reader should know where they are going

\section{Introduction}

% TODO: Write after Sections II-IX are complete
% Frame: "You know gauge fields, anomalies, Wilson lines, instantons, θ-terms.
%        When do these become 'topological' in a stronger sense?
%        This paper answers that question through Chern-Simons theory."
